% !TeX program = xelatex
\documentclass[a4paper,10pt]{article}

% ============================================================================
% កញ្ចប់បណ្ណាល័យ (Packages) – តាមលំនាំ NGS Preah Krasang (ls_plan.tex)
% ============================================================================
\usepackage[a4paper, left=0.95cm, right=0.95cm, top=0.85cm, bottom=0.92cm, footskip=0.52cm]{geometry}
\usepackage{amsmath}
\usepackage{unicode-math}
\usepackage{fontspec}
\usepackage{array, longtable, multirow, colortbl}
\usepackage[table]{xcolor}
\usepackage{enumitem}
\usepackage{pifont}
\usepackage{fancyhdr}
\usepackage{tikz}
\usetikzlibrary{calc, arrows.meta, angles, quotes, decorations.markings, babel}

% LaTeX 2025 + longtable v4.24 compatibility guard
\makeatletter
\let\@expl@@@mark@update@singlecol@structures@@\relax
\makeatother

% ============================================================================
% ការកំណត់ពុម្ពអក្សរ (Fonts Configuration - Cross-Platform Windows + Cloud fonts/)
% ============================================================================
\providecommand{\NGSFindFonts}{%
  \IfFileExists{C:/Windows/Fonts/cambria.ttc}{%
    \def\NGSSysPath{C:/Windows/Fonts/}%
    \def\NGSUserPath{C:/Users/KIT/AppData/Local/Microsoft/Windows/Fonts/}%
  }{%
    \IfFileExists{./fonts/KhmerOSsiemreap.ttf}{\def\NGSSysPath{./fonts/}\def\NGSUserPath{./fonts/}}{%
    \IfFileExists{../fonts/KhmerOSsiemreap.ttf}{\def\NGSSysPath{../fonts/}\def\NGSUserPath{../fonts/}}{%
    \IfFileExists{../../fonts/KhmerOSsiemreap.ttf}{\def\NGSSysPath{../../fonts/}\def\NGSUserPath{../../fonts/}}{%
    \IfFileExists{../../../fonts/KhmerOSsiemreap.ttf}{\def\NGSSysPath{../../../fonts/}\def\NGSUserPath{../../../fonts/}}{%
    \IfFileExists{../../../../fonts/KhmerOSsiemreap.ttf}{\def\NGSSysPath{../../../../fonts/}\def\NGSUserPath{../../../../fonts/}}{%
    \IfFileExists{../../../../../fonts/KhmerOSsiemreap.ttf}{\def\NGSSysPath{../../../../../fonts/}\def\NGSUserPath{../../../../../fonts/}}{%
    \IfFileExists{../../../../../../fonts/KhmerOSsiemreap.ttf}{\def\NGSSysPath{../../../../../../fonts/}\def\NGSUserPath{../../../../../../fonts/}}{%
      \def\NGSSysPath{../../../../../../../fonts/}\def\NGSUserPath{../../../../../../../fonts/}%
    }}}}}}}%
  }%
}
\NGSFindFonts

\setmathfont[
  Path={\NGSSysPath},
  FontIndex=1,
  Scale=1.02
]{cambria.ttc}

\defaultfontfeatures{Ligatures=TeX}

\setmainfont[
  Path={\NGSUserPath},
  Extension=.ttf,
  Script=Khmer,
  BoldFont={KhmerOSsiemreap},
  BoldFeatures={FakeBold=1.8},
  ItalicFont={KhmerOSsiemreap},
  ItalicFeatures={FakeSlant=0.25},
  BoldItalicFont={KhmerOSsiemreap},
  BoldItalicFeatures={FakeBold=1.8, FakeSlant=0.25}
]{KhmerOSsiemreap}

\newfontfamily\muol[
  Path={\NGSUserPath},
  Extension=.ttf,
  Script=Khmer
]{KhmerOSmuollight}

\newfontfamily\NagaBold[
  Path={\NGSUserPath},
  Extension=.ttf
]{AKbalthom Naga Bold}

\newfontfamily\CalibriFont[
  Path={\NGSSysPath},
  Extension=.ttf
]{calibri}

\newfontfamily\Tacteing[
  Path={\NGSUserPath}
]{TACTENG.TTF}

\newfontfamily\CourierBullet[
  Path={\NGSSysPath},
  Extension=.ttf
]{cour}
\newcommand{\wcirc}{{\CourierBullet\fontsize{9.5pt}{11pt}\selectfont o}}

% ក្បួនកាត់ពាក្យខ្មែរ (Khmer word wrapping)
\XeTeXlinebreaklocale "km"
\XeTeXlinebreakskip = 0pt plus 1pt minus 0.1pt

% ============================================================================
% ការកំណត់ពណ៌ និងទម្រង់ទំព័រ (Colors & Page Style)
% ============================================================================
\definecolor{stepgray}{HTML}{D9D9D9}
\definecolor{headgray}{HTML}{F2F2F2}
\definecolor{boxgray}{HTML}{D9D9D9}
\definecolor{darkred}{HTML}{C00000}
\definecolor{red5e}{HTML}{FF0000}
\definecolor{triblue}{HTML}{DAE8FC}
\definecolor{triborders}{HTML}{6C8EBF}
\definecolor{arrowblue}{HTML}{4472C4}
\definecolor{figblue}{HTML}{1F4E9E}
\definecolor{tealink}{HTML}{0F766E}
\definecolor{amberwarm}{HTML}{D97706}
\definecolor{pondcyan}{HTML}{DDF3FA}
\definecolor{pondstroke}{HTML}{0E7490}

\pagestyle{fancy}
\fancyhf{}
\renewcommand{\headrulewidth}{0pt}
\fancyfoot[R]{\CalibriFont\fontsize{11pt}{13pt}\selectfont\thepage}

\setlength{\parindent}{0pt}
\setlength{\parskip}{0pt}
\setlength{\tabcolsep}{3.8pt}

\newcolumntype{L}[1]{>{\raggedright\arraybackslash}p{#1}}

% ============================================================================
% ម៉ាក្រូជំនួយ (Helper Macros)
% ============================================================================
\newcommand{\fiveE}[1]{{\NagaBold\fontsize{9.5pt}{13pt}\selectfont\textcolor{red5e}{#1}}}
\newcommand{\timenote}[1]{\textcolor{darkred}{\textit{#1}}}
\newcommand{\redsub}[1]{\textcolor{darkred}{#1}}
\newcommand{\graybox}[1]{\colorbox{boxgray}{\strut\hspace{2pt}#1\hspace{2pt}}}
\newcommand{\ditem}{\tikz[baseline=-0.5ex]\fill (0,0.42ex) -- (0.42ex,0) -- (0,-0.42ex) -- (-0.42ex,0) -- cycle;\hspace{0.3em}}
\newcommand{\cnum}[1]{{\fontsize{10.5pt}{12pt}\selectfont\ding{\numexpr181+#1\relax}}\hspace{0.2em}}
\newcommand{\cfig}[1]{\par\vspace{4pt}{\centering #1\par}\vspace{4pt}}
\newcommand{\stepheader}[1]{\multicolumn{3}{|c|}{\cellcolor{stepgray}\muol #1} \\* \noalign{\hrule height 0.4pt}}
\newcommand{\tri}[1]{\ensuremath{\triangle #1}}
\newcommand{\ang}[1]{\ensuremath{\angle #1}}
\newcommand{\cg}{\ensuremath{\cong}}

% រចនាប័ទ្មគូសសម្គាល់ជ្រុងស្មើគ្នា (Tick mark styles - bold B&W print safe)
\tikzset{
  tk1/.style={postaction={decorate},decoration={markings,mark=at position 0.50 with {\draw[black, line width=0.85pt](0,-0.095)--(0,0.095);}}},
  tk2/.style={postaction={decorate},decoration={markings,
    mark=at position 0.45 with {\draw[black, line width=0.85pt](0,-0.095)--(0,0.095);},
    mark=at position 0.55 with {\draw[black, line width=0.85pt](0,-0.095)--(0,0.095);}}}
}

% ក្បាច់ក្រោមបាវចនា (Tacteing Font - Keyboard '3')
\newcommand{\khmflourish}{{\Tacteing\fontsize{36pt}{14pt}\selectfont\raisebox{-5pt}{3}}}

% ============================================================================
% ម៉ាក្រូគូសរូប TikZ (អក្សរពណ៌ខ្មៅ ១០០%, ទំហំសមាមាត្រ 3.8cm–6.8cm មិនបញ្ជ្រៀតគ្នា)
% every node/.style={font=\footnotesize, text=black} និងគម្លាត label >= 2.5mm
% ============================================================================
% \newcommand{\figExample}{%
% \begin{tikzpicture}[scale=1.0, every node/.style={font=\footnotesize, text=black}]
%   % ទទឹងរូបទោល 3.8cm–5.2cm (ក្នុង Col 2: 4.5cm–6.8cm) — ដាក់តែ ១ រូបក្នុង ១ ជួរដេកពីលើចុះក្រោម (ហាមដាក់ ២–៣ រូបទន្ទឹមគ្នា)
%   ...
% \end{tikzpicture}}

% ============================================================================
% ចាប់ផ្ដើមឯកសារ (Document Body — ចំនួនទំព័រ 4–8+ តាមខ្លឹមសារជាក់ស្ដែង, ពេញទំព័រ >=90%)
% ============================================================================
\begin{document}
\fontsize{10pt}{15.2pt}\selectfont

% -------------------- ក្បាលកិច្ចតែងការ (Header) --------------------
\begin{center}
  {\muol\fontsize{11pt}{16pt}\selectfont ព្រះរាជាណាចក្រកម្ពុជា}\\[2pt]
  {\muol\fontsize{10pt}{15pt}\selectfont ជាតិ សាសនា ព្រះមហាក្សត្រ}\\[1pt]
  \khmflourish
\end{center}

\vspace{-15pt}
{\muol\fontsize{9.5pt}{13.2pt}\selectfont
មន្ទីរអប់រំ យុវជន និងកីឡា ខេត្តកំពង់ធំ\\
ការិយាល័យអប់រំ យុវជន និងកីឡា ស្រុកកំពង់ស្វាយ\\
អនុវិទ្យាល័យព្រះក្រសាំង}

\vspace{2pt}
\begin{center}
  {\muol\fontsize{10pt}{15pt}\selectfont កិច្ចតែងការបង្រៀន}
\end{center}
\vspace{1pt}

% -------------------- ព័ត៌មានទូទៅ (Metadata) --------------------
{\renewcommand{\arraystretch}{1.06}
\hspace*{4.6cm}%
\begin{tabular}{@{}l l l@{}}
  - \hspace{0.35cm} មុខវិជ្ជា & : & គណិតវិទ្យា ថ្នាក់ទី<<ថ្នាក់>> \\
  - \hspace{0.35cm} ចំណងជើងមេរៀន & : & <<មេរៀន>> \\
  - \hspace{0.35cm} ចំណងជើងរង & : & <<ចំណងជើងរង>> \\
  - \hspace{0.35cm} រយៈពេល & : & ១ ម៉ោងសិក្សា (៥០នាទី) \\
  - \hspace{0.35cm} បង្រៀនដោយ & : & ហុក ប៊ុនថែន \\
  - \hspace{0.35cm} កាលបរិច្ឆេទ & : & \makebox[28mm]{\dotfill}\ ២០២៦ \\
\end{tabular}

\vspace{2pt}

% -------------------- I. វត្ថុបំណង --------------------
\textbf{I.} {\muol វត្ថុបំណង}\\[2pt]
\hspace*{0.55cm}%
\begin{tabular}{@{}l l l@{}}
  \wcirc \hspace{0.25cm} វិជ្ជាសម្បទា & : & <<វិជ្ជាសម្បទា>> \\
  \wcirc \hspace{0.25cm} បំណិនសម្បទា & : & <<បំណិនសម្បទា>> \\
  \wcirc \hspace{0.25cm} ចរិយាសម្បទា & : & <<ចរិយាសម្បទា>> \\
\end{tabular}

\vspace{2pt}

% -------------------- II. សម្ភារៈឧបទេស --------------------
\textbf{II.} {\muol សម្ភារៈឧបទេស}\\[2pt]
\hspace*{0.55cm}$\bullet$ \hspace{0.35cm} <<សម្ភារៈឧបទេស>>

\vspace{2pt}

% -------------------- III. វិធីសាស្រ្តបង្រៀន --------------------
\textbf{III.} {\muol វិធីសាស្រ្តបង្រៀន}\\[2pt]
\hspace*{0.55cm}%
\begin{tabular}{@{}l l l@{}}
  $\bullet$ \hspace{0.35cm} គោលវិធី & : \hspace{0.7cm} & សិស្សមជ្ឈមណ្ឌល \\
  $\bullet$ \hspace{0.35cm} យុទ្ធវិធីបង្រៀន & : & ស្ថាបនានិយម គម្រូ 5Es (<<សកម្មភាពសង្ខេប>>) \\
\end{tabular}
}

\vspace{2pt}

% -------------------- IV. ដំណើរការបង្រៀន --------------------
\textbf{IV.} {\muol ដំណើរការបង្រៀន}

\vspace{-3pt}

\renewcommand{\arraystretch}{1.15}
\begin{longtable}{|L{5.50cm}|L{7.90cm}|L{4.80cm}|}
\hline
\rowcolor{headgray}
\multicolumn{1}{|c|}{\muol សកម្មភាពគ្រូ} &
\multicolumn{1}{c|}{\muol ខ្លឹមសារមេរៀន} &
\multicolumn{1}{c|}{\muol សកម្មភាពសិស្ស} \\
\hline
\endfirsthead

\hline
\endhead

\hline
\endfoot

\hline
\endlastfoot

% ==================== ជំហានទី១ (2min) ====================
\stepheader{ជំហានទី១ (រដ្ឋបាលថ្នាក់)}
- ពិនិត្យអវត្តមាន និងអនាម័យថ្នាក់រៀន
&
\hspace*{0.3cm}- រដ្ឋបាលថ្នាក់
&
- រាយការណ៍អវត្តមាន \\[3pt]
\hline

% ==================== ជំហានទី២ (9min: ≤3 + 3 + E1 3) ====================
\stepheader{ជំហានទី២ (កិច្ចការផ្ទះ + ទំនាក់ទំនងមេរៀនថ្មី)}
% ច្បាប់បំបែកជួរ longtable: រាល់ E1..E5 ត្រូវបំបែកជាជួររងតូចៗ (2–5 បន្ទាត់/ជួរ) បញ្ចប់ដោយ \\[3pt] (គ្មាន \hline កណ្ដាលជំហាន)
% រូបភាព TikZ នីមួយៗដាក់ក្នុងជួររងដាច់ដោយឡែក (& \cfig{...} & \\[3pt]) ដើម្បីឱ្យទំព័រ 1..N-1 ពេញ >=90%
% គំរូជួររងអក្សរចុះបន្ទាត់ (Teacher-Friendly Vertical Structure — ហាមសរសេរអក្សរតរហូត):
% \fiveE{E1 - Engage} \timenote{(3 min)}\par\vspace{2pt}
% $\bullet$ គ្រូបង្ហាញស្ថានភាព...\par\vspace{2pt}
% \hspace*{0.3em}-- សំណួរ៖ «...»
% &
% \textbf{១. បេសកកម្ម E1៖ ...}\par\vspace{2pt}
% ...
% &
% $\circ$ សិស្សសង្កេត...\par\vspace{2pt}
% $\circ$ សិស្សទស្សន៍ទាយ... \\[3pt]
% & \cfig{\figExample} & \\[3pt]
\hline

% ==================== ជំហានទី៣ (17min: E2 11 + E3 6) ====================
\stepheader{ជំហានទី៣ (ខ្លឹមសារមេរៀនថ្មី)}
% បំបែក E2 (11 min) និង E3 (6 min) ជាជួររងតូចៗ 2–5 បន្ទាត់ + ជួររងរូបភាពដាច់ដោយឡែក \\[3pt]
% គ្រប់ជំហានគណិតវិទ្យា សម្មតិកម្ម/សន្និដ្ឋាន/សម្រាយបញ្ជាក់ និងចម្លើយសិស្ស ត្រូវចុះបន្ទាត់ថ្មីជានិច្ច
& & \\[3pt]
\hline

% ==================== ជំហានទី៤ (20min: E4 10 + E5 10) ====================
\stepheader{ជំហានទី៤ (ពង្រីកចំណេះដឹង)}
% បំបែក E4 (10 min) និង E5 (10 min) ជាជួររងតូចៗ + ជួររងរូបភាពធំច្បាស់ 3.8–6.8 cm \\[3pt]
& & \\[3pt]
\hline

% ==================== ជំហានទី៥ (≤2min) ====================
\stepheader{ជំហានទី៥ (កិច្ចការផ្ទះ + បណ្ដាំផ្ញើ)}
\timenote{ចំណាយពេលឱ្យតិចបំផុតកុំឱ្យលើស 2min}\par\vspace{2pt}
$\bullet$ គ្រូដាក់កិច្ចការផ្ទះ...
&
\centerline{\muol កិច្ចការផ្ទះ}\par\vspace{2pt}
% ១ បន្ទាត់ = ១ លំហាត់/សំណួរ
&
$\circ$ សិស្សកត់ត្រា \\[3pt]
\end{longtable}

% -------------------- កាលបរិច្ឆេទ និងហត្ថលេខា --------------------
\vspace{2pt}
\hfill
\begin{minipage}{9.8cm}
  {\centering
    ថ្ងៃ\makebox[16mm]{\dotfill}\ ខែ\makebox[18mm]{\dotfill}\ ឆ្នាំមមី អដ្ឋស័ក ពុទ្ធសករាជ ២៥៧០\\[2pt]
    ត្រូវនឹងថ្ងៃទី\makebox[10mm]{\dotfill}\ ខែ\makebox[18mm]{\dotfill}\ ឆ្នាំ២០២៦\\[3pt]
    {\muol គ្រូបង្រៀន}\par}
  \vspace{32pt}
  \hfill{\muol ហុក ប៊ុនថែន}\hspace*{1.4cm}
\end{minipage}

\end{document}
